\section{Гауссовский случайный вектор.}
\subsection{Теоретическая справка.}
\begin{definition}
    Случайный вектор $\vec{X}$ со значениями в $\R^d$ будем называть гауссовским, если для всякого вектора $\vec{t} \in \R^d$ случайная величина $\langle \vec{X}, \vec{t} \rangle$~---~гауссовская.
\end{definition}
\begin{note}
    Очевидно, что всякая компонента гауссовского вектора -- гауссовская случайная величина. \\
    При этом, неверно, что если выбрать $n$ гауссовских случайных величин и составить из них вектор, то он будет гауссовским. \\
    Впрочем, для независимых гауссовских случайных величин это утверждение справедливо.
\end{note}
\begin{proposition}
    Если $\vec{X}$~---~гауссовский, то для всякой прямоугольной матрицы подходящего размера вектор $\vec{Y} = A\vec{X} + \vec{b}$~---~гауссовский.
\end{proposition}
\begin{proposition}
    Гауссовский случайный вектор полностью определяется своим математическим ожиданием $\vec{m} = \Ex \vec{X}$ и ковариационной матрицей $\D \vec{X} = R$.
    Многомерное гауссовское распределение обозначается как $N(\vec{m}, R)$.
\end{proposition}
\begin{proposition}
    Если $\vec{X} \sim N(\vec{m}, R)$, то $\vec{Y} = A\vec{X} + \vec{b}$ распределён как $\vec{Y} \sim N(A\vec{m} + \vec{b}, ARA^T)$.
\end{proposition}
\begin{proposition}
    Если $\vec{Y} \sim N(\vec{m}, R)$ (размерности $d$), то существует такое линейное преобразование $A: \R^d \to \R^d$, что $\vec{Y} = A\vec{X} + \vec{b}$, где $\vec{X}$ распределён как $N(\vec{0}, I)$.
\end{proposition}
\begin{lemma}
    Если $X \sim N(a_1, \sigma_1^2)$ и $Y \sim N(a_2, \sigma_2^2)$~----~независимые случайные величины, то $X + Y \sim N(a_1 + a_2, \sigma_1^2 + \sigma_2^2)$.
\end{lemma}
\begin{theoremdefinition}
    Гауссовским случайным вектором будем назвыать случайный вектор с характеристической функцией $e^{i\vec{m}^T\vec{t} - \frac{1}{2}\vec{t}^{T}R\vec{t}}$
\end{theoremdefinition}
\begin{proposition}
    Для компонент гауссовского случайного вектора независимость и некоррелированность -- одно и то же.
\end{proposition}
\begin{proposition}
    Пусть $\vec{X} \sim N(\vec{m}, R), \vec{X} \in \R^d$ и $R$~---~невырожденная матрица.
    Тогда существует плотность равная
    \[
        p_{\vec{X}}(\vec{x}) = \dfrac{1}{\sqrt {2\pi}^d \sqrt {|\det R|}}\exp\biggr(-\dfrac{1}{2}(\vec{x} - \vec{m})^TR^{-1}(\vec{x} - \vec{m})\biggr).
    \]
\end{proposition}